\documentclass[11pt]{amsart}

\usepackage{amsmath,amssymb,mathtools}
\usepackage{microtype}
\usepackage[hidelinks]{hyperref}

\hypersetup{
  pdftitle={Factorial-residue hypergraphs and an auxiliary problem of Erdos},
  pdfauthor={Aakash Gurung},
  pdfsubject={Weighted factorial-residue hypergraphs and rough integers},
  pdfkeywords={factorials, edge-colored hypergraphs, rainbow edge covers, sieve methods, rough numbers, covering congruences}
}

\newtheorem{theorem}{Theorem}[section]
\newtheorem{proposition}[theorem]{Proposition}
\newtheorem{lemma}[theorem]{Lemma}
\theoremstyle{definition}
\newtheorem{definition}[theorem]{Definition}
\theoremstyle{remark}
\newtheorem{remark}[theorem]{Remark}

\title[Factorial-residue hypergraphs]
{Factorial-residue hypergraphs and an auxiliary problem of Erd\H{o}s}
\author{Aakash Gurung}
\address{Department of Mathematics, University of Alabama,
Tuscaloosa, AL 35487, USA}
\date{}

\subjclass[2020]{11N36, 11B65, 11A41, 05C65}
\keywords{factorials, edge-colored hypergraphs, rainbow edge covers, sieve methods, rough numbers, covering congruences}

\begin{document}

\begin{abstract}
For $L\geq2$, we associate a weighted edge-colored multihypergraph with the
factorial residues modulo primes $q>L$.  We prove that the minimum weight
$\kappa(L)$ of a rainbow edge cover is at most $(1-\delta)L\log L$ for some
absolute $\delta>0$ and all sufficiently large $L$.  This minimum is the
logarithm of the least modulus arising from the corresponding system of
covering congruences.  As an application, we prove that there are absolute
constants $\eta>0$ and $L_0$ such that, for every $L\geq L_0$, at least
\[
\frac{((L+1)!)^\eta}{100\log L}
\]
integers $n\in(2L!,(L+1)!]$ satisfy $P^-(n)>L$ and have $n-k!$ composite
for every $1\leq k\leq L$.  This gives a quantitative solution of an
auxiliary problem of Erd\H{o}s.
\end{abstract}

\maketitle

\section{Introduction}

Guy \cite[Problem A2]{Guy} records the following question of Erd\H{o}s: are
there infinitely many primes $p$ such that $p-k!$ is composite whenever
$1\leq k!<p$?  The primes $101$ and $211$ are examples.  This question is
also listed as Erd\H{o}s Problem 1059 and remains open \cite{Bloom}.

Alongside Problem 1059, Erd\H{o}s suggested the following easier auxiliary
problem \cite[Problem A2]{Guy}: prove that there are infinitely many pairs
$(L,n)$ with
\[
L!<n\leq(L+1)!,
\]
such that every prime factor of $n$ exceeds $L$ and $n-k!$ is composite for
every $1\leq k\leq L$.

We prove this auxiliary problem in a quantitative form.  For an integer
$n>1$, let $P^-(n)$ and $P^+(n)$ denote its least and greatest prime factors,
respectively.  We call $n$ \emph{$L$-rough} if $P^-(n)>L$.

The proof begins with a weighted edge-colored multihypergraph on the vertex
set $\{2,\ldots,L\}$.  Whenever several factorials have the same nonzero
residue modulo a prime $q>L$, their indices form an edge of color $q$ and
weight $\log q$.  Let $\kappa(L)$ be the minimum weight of a rainbow edge
cover, where rainbow means that no two selected edges have the same color.
The precise definition is given in Section~2.

\begin{theorem}\label{thm:covercost}
There are absolute constants $\delta>0$ and $L_0$ such that
\[
\kappa(L)\leq(1-\delta)L\log L
\]
for every integer $L\geq L_0$.
\end{theorem}

The quantity $\kappa(L)$ has an exact arithmetic meaning: it is the
logarithm of the least modulus arising from a system of covering
congruences of the prescribed form.  Thus, Theorem~\ref{thm:covercost}
gives a fixed proportional saving at the $L\log L$ scale.

\begin{theorem}\label{thm:main}
There are absolute constants $\eta>0$ and $L_1$ such that, for every integer
$L\geq L_1$, if $X=(L+1)!$, then at least
\[
\frac{X^\eta}{100\log L}
\]
integers $n\in(2L!,X]$ satisfy $P^-(n)>L$ and have $n-k!$ composite for
every $1\leq k\leq L$.
\end{theorem}

In particular, solutions of the auxiliary problem exist for every
sufficiently large $L$.  They lie in the smaller interval
$(2L!,(L+1)!]$.  The theorem does not say that any of these integers is
prime, so it does not solve Erd\H{o}s Problem 1059.

Factorial residues modulo a fixed prime have been studied from the viewpoints
of value sets and distribution; see, for example,
\cite{BanksLucaShparlinskiStichtenoth,GaraevLucaShparlinski,
KlurmanMunsch,Hu}.  Here the prime varies with the edge.  An auxiliary
conflict graph separates the combinatorial selection of a rainbow subfamily
from the arithmetic problem of producing sufficiently many low-weight
edges.

For the unconditional estimate, we use collisions of gap two.  The basic
identity is
\begin{equation}\label{eq:keyidentity}
(i+2)!-i!=i!\bigl(i^2+3i+1\bigr).
\end{equation}
If a prime $q>i+2$ divides $i^2+3i+1$, then one fiber modulo $q$ contains
both $i$ and $i+2$.  A product argument and a one-dimensional upper-bound
sieve give linearly many candidate fibers with
\[
L<q\ll L^{2-\alpha}
\]
for a fixed $\alpha>0$.  An independent-set argument selects a constant
proportion having distinct colors and disjoint designated pairs.  The
remaining vertices are covered separately by primes of size $O(L\log L)$.
The paired fibers save a fixed multiple of $L\log L$, which proves
Theorem~\ref{thm:covercost}.

We then prove a general transference theorem.  Any rainbow edge cover whose
modulus is at most $((L+1)!)^{1-\eta}$ produces many integers in one CRT
progression.  Counting reduced residue classes modulo the primorial
$\prod_{p\leq L}p$ yields the $L$-rough integers.  The restrictions
$q<L!$ and $n>2L!$ make every selected divisor of $n-k!$ proper, and the
case $k=1$ follows from parity.  This gives Theorem~\ref{thm:main}.

\section{The factorial-residue hypergraph}

Let $L\geq2$ and put
\[
V_L=\{2,3,\ldots,L\}.
\]

\begin{definition}\label{def:fibres}
For a prime $L<q<L!$ and a nonzero residue $a\pmod q$, let $e_{q,a}$ be a
labeled hyperedge with incident vertex set
\[
E_{q,a}=\{k\in V_L:k!\equiv a\pmod q\}.
\]
After empty edges are omitted, these labeled edges form a weighted
edge-colored multihypergraph $\mathcal H_L$.  The edge $e_{q,a}$ has color
$q$ and weight $\log q$.
\end{definition}

Since $q>L$, every $k!$ with $k\in V_L$ is nonzero modulo $q$.  Therefore,
for each fixed $q$, the fibers $E_{q,a}$ partition $V_L$, with empty fibers
omitted.

\begin{definition}\label{def:cover}
A \emph{rainbow edge cover} of $\mathcal H_L$ is a finite collection
\[
\mathcal C=\{e_{q,a_q}:q\in Q\}
\]
with distinct primes $L<q<L!$, one residue $a_q\in(\mathbb Z/q\mathbb Z)^\times$
for each $q\in Q$, nonempty fibers $E_{q,a_q}$, and
\[
V_L\subseteq\bigcup_{q\in Q}E_{q,a_q}.
\]
Its weight is
\[
w(\mathcal C)=\sum_{q\in Q}\log q.
\]
We define
\[
\kappa(L)=\inf_{\mathcal C}w(\mathcal C),
\]
where the infimum is over all rainbow edge covers, and put
$\kappa(L)=+\infty$ if none exists.
\end{definition}

The rainbow condition is necessary because choosing two different fibers of
the same color would impose incompatible residues modulo the same prime.

\begin{proposition}[Collision criterion]\label{prop:criterion}
Let $2\leq i<j\leq L$ and let $q>L$ be prime.  Then
\[
\begin{aligned}
i!\equiv j!\pmod q
&\quad\Longleftrightarrow\quad
\prod_{r=i+1}^{j}r\equiv1\pmod q\\
&\quad\Longleftrightarrow\quad
q\mid\left(\prod_{r=i+1}^{j}r-1\right).
\end{aligned}
\]
\end{proposition}

\begin{proof}
We have
\[
j!=i!\prod_{r=i+1}^{j}r.
\]
Also, $q>L\geq i$, so $q\nmid i!$.  We may therefore cancel $i!$ modulo
$q$.  The three conditions are now equivalent.
\end{proof}

The next proposition proves that the scale $L\log L$ is the natural
singleton baseline.

\begin{proposition}\label{prop:baseline}
For every sufficiently large $L$,
\[
\kappa(L)\leq(L-1)\log(20L\log L)
=(1+o(1))L\log L.
\]
In particular, a rainbow edge cover exists, and the infimum in
Definition~\ref{def:cover} is a minimum.
\end{proposition}

\begin{proof}
The prime number theorem gives
\[
\pi(20L\log L)-\pi(L)>2L
\]
for all sufficiently large $L$.  Choose distinct primes
$q_k\in(L,20L\log L]$ for $k\in V_L$.  We also have
$20L\log L<L!$ for large $L$.  The edges $e_{q_k,k!}$ form a rainbow edge
cover, and their total weight is at most
\[
(L-1)\log(20L\log L).
\]
There are only finitely many primes between $L$ and $L!$, and each has only
finitely many fibers.  Hence there are only finitely many rainbow edge covers.
The nonempty finite set of their weights has a minimum.
\end{proof}

We next give the exact connection with the Chinese remainder theorem.  For
sufficiently large $L$, let $M_{\min}(L)$ be the least squarefree product
$M$ of primes $L<q<L!$ for which there exists
$A\in(\mathbb Z/M\mathbb Z)^\times$ satisfying both of the following
conditions:
\begin{enumerate}
\item for every prime $q\mid M$, there is a $k\in V_L$ such that
$A\equiv k!\pmod q$;
\item for every $k\in V_L$, there is a prime $q\mid M$ such that
$A\equiv k!\pmod q$.
\end{enumerate}

\begin{proposition}[Equivalence with covering congruences]
\label{prop:crt-equivalence}
For every sufficiently large $L$,
\[
\kappa(L)=\log M_{\min}(L).
\]
More precisely, a rainbow edge cover $\mathcal C$ gives a modulus
\[
M(\mathcal C)=\prod_{q\in Q}q
\]
with
\[
\log M(\mathcal C)=w(\mathcal C),
\]
and every pair $(M,A)$ in the definition of $M_{\min}(L)$ gives a rainbow
edge cover of weight $\log M$.
\end{proposition}

\begin{proof}
Let $\mathcal C=\{e_{q,a_q}:q\in Q\}$ be a rainbow edge cover.  The primes
in $Q$ are distinct, so the Chinese remainder theorem gives a unique residue class
$A\pmod {M(\mathcal C)}$ satisfying
\[
A\equiv a_q\pmod q\qquad(q\in Q).
\]
Every $a_q$ is nonzero, so $(A,M(\mathcal C))=1$.  If $k\in E_{q,a_q}$,
then $A\equiv k!\pmod q$.  Thus $(M(\mathcal C),A)$ satisfies the defining
conditions for $M_{\min}(L)$.  Also,
\[
\log M(\mathcal C)=\sum_{q\in Q}\log q
=w(\mathcal C).
\]

Conversely, let $(M,A)$ satisfy the defining conditions for $M_{\min}(L)$.
For each prime $q\mid M$, choose the edge $e_{q,A\bmod q}$.  The residue is
nonzero because $(A,M)=1$, and the second condition says that these edges
cover $V_L$.  They use at most one edge of each color and have total weight
\[
\sum_{q\mid M}\log q=\log M.
\]
Taking the two minima proves the proposition.
\end{proof}

For the selection argument, let $\mathcal F$ be a finite family of triples
\[
F=(q,a,B_F),
\qquad
\varnothing\neq B_F\subseteq E_{q,a}.
\]
The set $B_F$ is the \emph{designated subset} of $F$.  Define the auxiliary
conflict graph $G_{\mathcal F}$ on vertex set $\mathcal F$ by joining
distinct triples $F=(q,a,B_F)$ and $F'=(q',a',B_{F'})$ whenever $q=q'$ or
\[
B_F\cap B_{F'}\neq\varnothing.
\]

\begin{proposition}[Weighted independent-set lemma]
\label{prop:conflict-selection}
Suppose that the conflict graph of $\mathcal F$ has maximum degree at most
$D$.  If $w(F)\geq0$ for every $F\in\mathcal F$, then there is a subfamily
$\mathcal J\subseteq\mathcal F$ with distinct prime colors and pairwise
disjoint designated subsets such that
\[
\sum_{F\in\mathcal J}w(F)
\geq\frac1{D+1}\sum_{F\in\mathcal F}w(F).
\]
In the special case $w(F)=1$ for every $F$, one may choose $\mathcal J$ so
that
\[
|\mathcal J|\geq\frac{|\mathcal F|}{D+1}.
\]
\end{proposition}

\begin{proof}
A finite graph of maximum degree at most $D$ has a proper coloring with at
most $D+1$ colors: order the vertices arbitrarily and color them greedily.
The color classes are independent sets and partition $\mathcal F$.  One of
them has weight at least the average of their weights.  Taking this color
class as $\mathcal J$ proves the weighted inequality.  Taking unit weights
gives the last assertion.  Independence gives distinct prime colors and
pairwise disjoint designated subsets by the definition of the conflict graph.
\end{proof}

Relative to the leading weight $|B_F|\log L$ of $|B_F|$ singleton edges, an
edge of color $q$ yields a power saving only when
$q\leq L^{|B_F|-\varepsilon}$ for some $\varepsilon>0$.  In the quadratic
construction $|B_F|=2$, which explains the range $q\ll L^{2-\alpha}$.

\section{Prime divisors of a quadratic}

We first estimate the frequency with which $i^2+3i+1$ has a sufficiently
large prime divisor.  Throughout this section, put
\[
f(t)=t^2+3t+1.
\]

\subsection{Local root counts}

We also let $\chi_5$ be the primitive quadratic character modulo $5$, and we
extend it by setting $\chi_5(5)=0$.  We will use the prime number theorem, the
Siegel--Walfisz theorem for the fixed modulus $5$, and Mertens's estimates in
the forms stated in
\cite[Theorems 3.4, 8.1 and 12.1]{Koukoulopoulos}.  Let
\[
\Theta_{\chi_5}(x)=\sum_{p\leq x}\chi_5(p)\log p.
\]
For $a\in(\mathbb Z/5\mathbb Z)^\times$, the fixed-modulus Siegel--Walfisz
estimate for $\pi(x;5,a)$ and partial summation give
$\vartheta(x;5,a)=x/4+O(xe^{-c\sqrt{\log x}})$ for some $c>0$.  Since the
sum of $\chi_5(a)$ over these four classes is zero, summing with weights
$\chi_5(a)$ gives
$\Theta_{\chi_5}(x)\ll xe^{-c\sqrt{\log x}}$.
Partial summation therefore gives
\begin{equation}\label{eq:chi}
\begin{aligned}
\sum_{p\leq x}\chi_5(p)\frac{\log p}{p-1}
&=\frac{\Theta_{\chi_5}(x)}{x-1}
 +\int_2^x\frac{\Theta_{\chi_5}(t)}{(t-1)^2}\,dt+O(1)\\
&=O(1).
\end{aligned}
\end{equation}
Indeed, the integrand is $O(e^{-c\sqrt{\log t}}/t)$, which is integrable on
$[2,\infty)$.  The same argument, now using the weight $1/(t\log t)$,
gives
\[
\sum_{p\leq x}\frac{\chi_5(p)}p
=\frac{\Theta_{\chi_5}(x)}{x\log x}
 +\int_2^x\Theta_{\chi_5}(t)
 \frac{\log t+1}{t^2(\log t)^2}\,dt+O(1)
=O(1).
\]
For $p\neq5$, we have
$\mathbf 1_{\chi_5(p)=1}=(1+\chi_5(p))/2$.  Mertens's estimate now gives
\begin{equation}\label{eq:split}
\sum_{\substack{p\leq x\\ \chi_5(p)=1}}\frac1p
=\frac12\log\log x+O(1).
\end{equation}

The identity
\begin{equation}\label{eq:discriminant}
(2t+3)^2=4f(t)+5
\end{equation}
shows that, for $p\neq2,5$, the number of roots of $f$ modulo $p$ is
$1+\left(\frac5p\right)$.  Since $5\equiv1\pmod4$, quadratic reciprocity
gives
\[
\left(\frac5p\right)=\left(\frac p5\right)=\chi_5(p).
\]
Thus,
\begin{equation}\label{eq:rho}
\rho(p)=1+\chi_5(p)\in\{0,2\}.
\end{equation}
The derivative is $f'(t)=2t+3$.  If a root of $f$ modulo $p$ were not simple,
then $f(t)\equiv f'(t)\equiv0\pmod p$, and \eqref{eq:discriminant} would give
$p\mid5$.  Hence, every root under consideration is simple.  Hensel's lemma
then gives exactly one lift of each root to every modulus $p^a$.  There is no
root modulo $2$.  Modulo $5$ there is one root, but it does not lift to a root
modulo $25$, because
\begin{equation}\label{eq:five}
f(1+5t)=5(1+5t+5t^2).
\end{equation}

\subsection{Large prime factors}

\begin{lemma}\label{lem:largeprime}
For all sufficiently large $N$,
\[
\#\left\{1\leq i\leq N:P^+(f(i))>\frac18 i\log i\right\}\geq\frac N5.
\]
\end{lemma}

\begin{proof}
Let $E_N$ denote the set of indices for which all prime factors of $f(i)$
are small, and define
\[
E_N=\left\{i:\frac{N}{\log N}<i\leq N,\ P^+(f(i))\leq i\right\},
\qquad
Q_N=\prod_{i\in E_N}f(i).
\]
Every prime divisor of $Q_N$ is at most $N$.  For $p\neq2,5$, the root count
modulo prime powers gives
\begin{align*}
\operatorname{ord}_p\left(\prod_{i\leq N}f(i)\right)
&=\sum_{a\geq1}\#\{i\leq N:p^a\mid f(i)\}.
\end{align*}
For every $a\geq1$, there are $\rho(p)$ root classes modulo $p^a$, and hence
\[
\#\{i\leq N:p^a\mid f(i)\}
\leq \frac{N\rho(p)}{p^a}+\rho(p).
\]
Also, $p^a\mid f(i)$ is possible only when $p^a\leq N^2+3N+1$.
Consequently, only $O(\log N/\log p)$ values of $a$ occur, and summing the
geometric series gives
\[
\operatorname{ord}_p\left(\prod_{i\leq N}f(i)\right)
\leq
\frac{N\rho(p)}{p-1}
+O\left(\rho(p)\frac{\log N}{\log p}\right).
\]
The prime $2$ contributes nothing.  Equation \eqref{eq:five} shows that
$5\mid f(i)$ precisely when $i\equiv1\pmod5$, and that in this case
$5^2\nmid f(i)$.  Thus, the contribution of $p=5$ to the logarithm below is
$O(N)$.  Moreover,
\[
\log N\sum_{p\leq N}\rho(p)\ll \log N\,\pi(N)\ll N.
\]
Therefore, Mertens's estimate together with \eqref{eq:chi} gives
\begin{align*}
\log Q_N
&\leq
N\sum_{p\leq N}\frac{(1+\chi_5(p))\log p}{p-1}
+O\left(\log N\sum_{p\leq N}\rho(p)\right)+O(N)\\
&=N\log N+O(N).
\end{align*}
For the corresponding lower bound, uniformly for
$N/\log N<i\leq N$, we have
\[
\log f(i)=2\log N+O(\log\log N).
\]
It follows that
\[
\log Q_N
\geq |E_N|\bigl(2\log N-O(\log\log N)\bigr).
\]
Comparing this inequality with the upper estimate for $\log Q_N$, we obtain
\[
|E_N|\leq\left(\frac12+o(1)\right)N.
\]
The interval $(N/\log N,N]$ contains $N-o(N)$ integers.  Thus, at least
$(1/2-o(1))N$ integers in this interval satisfy $P^+(f(i))>i$.  Put
\[
\mathcal S_N=
\left\{i:\frac{N}{\log N}<i\leq N,\ P^+(f(i))>i\right\}.
\]
We have proved that $|\mathcal S_N|\geq(1/2-o(1))N$.  It remains to
strengthen the lower bound for these prime factors.

Discard from $\mathcal S_N$ the indices for which
\[
P^+(f(i))\leq\frac18N\log N.
\]
For every discarded index, the prime $q=P^+(f(i))$ exceeds $i$.  Hence, the
relevant indices lie in $[1,q-1]$ and are distinct modulo $q$.  Since $f$
has at most two roots modulo $q$, a fixed prime $q$ can occur for at most two
indices.  Therefore,
the number of discarded indices is at most
\[
2\pi\left(\frac18N\log N\right)
=\left(\frac14+o(1)\right)N.
\]
We started with at least $(1/2-o(1))N$ indices and discarded at most
$(1/4+o(1))N$ of them.  Thus, at least $(1/4-o(1))N$ indices remain.  For
each remaining index,
\[
P^+(f(i))>\frac18N\log N\geq\frac18i\log i.
\]
Since $1/4-o(1)>1/5$ for all sufficiently large $N$, at least $N/5$ indices
satisfy the required inequality.  This proves the lemma.
\end{proof}

\section{A sieve estimate for the complementary factor}

A large prime divisor of $f(i)$ is useful only if its complementary factor
is not too small, since otherwise the prime is too costly for the covering
modulus.  Write $f(i)=mq$, where $q>L$ is prime.  We show that $m$ is not too
small for most of the indices under consideration.

\subsection{Sieve setup}

For $0<\alpha\leq1/2$, let
\[
E_\alpha(L)
=
\#\left\{1\leq i\leq L:
 f(i)=mq\text{ for some }m\leq L^\alpha
 \text{ and prime }q>L
\right\}.
\]
Thus, $E_\alpha(L)$ counts the indices for which a prime factor $q>L$ has a
complementary factor $m\leq L^\alpha$.  The next lemma says that this bad set
is small when $\alpha$ is small.

\begin{lemma}\label{lem:smallcofactor}
There is an absolute constant $C_0\geq1$ such that, uniformly for $0<\alpha\leq1/2$ and all sufficiently large $L$,
\[
E_\alpha(L)
\leq
C_0\left(
\alpha L+\frac{L}{\log L}
+L^{(1+\alpha)/2}(\log L)^2
\right).
\]
The lower threshold for $L$ may be chosen independently of $\alpha$.
\end{lemma}

\begin{proof}
We begin by counting the roots of $f$ modulo $m$.  For $m\geq1$, write
\[
R(m)=\#\{a\pmod m:f(a)\equiv0\pmod m\}.
\]
The prime-power values from the previous section are
\[
R(2^e)=0\quad(e\geq1),
\qquad
R(5)=1,
\qquad
R(5^e)=0\quad(e\geq2),
\]
and, for $p\neq2,5$,
\[
R(p^e)=1+\chi_5(p)\quad(e\geq1).
\]
The Chinese remainder theorem therefore gives
\begin{equation}\label{eq:R}
R(m)=
\begin{cases}
2^{\omega(m/(m,5))},
&\begin{array}{l}
2\nmid m,\ 25\nmid m,\ \text{and}\\
\chi_5(p)=1\text{ for every }p\mid m,\ p\neq5,
\end{array}\\[3mm]
0,&\text{otherwise}.
\end{cases}
\end{equation}

Now fix $m\leq L^\alpha$ and a root $a\pmod m$, represented by
$0\leq a<m$.  Every integer $i\equiv a\pmod m$ can be written as
$i=a+mt$.  Accordingly, let
\[
i=a+mt,
\qquad
T_{a,m}=\{t\in\mathbb Z:1\leq a+mt\leq L\},
\qquad
H=|T_{a,m}|.
\]
The length of this interval of $t$ values is
\[
H=\frac Lm+O(1),
\qquad
H\asymp\frac Lm,
\qquad
H\geq L^{1/2}-1.
\]
Since $m\mid f(a)$, we can divide $f(a+mt)$ by $m$ and obtain the integral
polynomial
\[
    G_{a,m}(t)=\frac{f(a+mt)}m
  =mt^2+(2a+3)t+\frac{f(a)}m
  .\]
This polynomial has discriminant $5$.  It is also primitive.  Indeed, if a
prime divided all three coefficients, then its square would divide
$(2a+3)^2-4f(a)=5$, which is impossible.

If $\ell\nmid10m$ is prime, then the affine map $t\mapsto a+mt$ is a
bijection modulo $\ell$.  Consequently,
\[
\#\{t\pmod\ell:G_{a,m}(t)\equiv0\pmod\ell\}
=1+\chi_5(\ell).
\]
Let
  \[
  \mathcal P_0=\{\ell:\ell\text{ is prime and }\ell\nmid10\},
  \qquad
  \mathcal P_m=\{\ell\in\mathcal P_0:\ell\nmid m\},
  \]
and extend $\rho$ multiplicatively to squarefree integers supported on
$\mathcal P_m$, with $\rho(d)=0$ otherwise.  For supported squarefree $d$,
counting the root classes inside $T_{a,m}$ gives
\begin{equation}\label{eq:localcount}
\#\{t\in T_{a,m}:d\mid G_{a,m}(t)\}
=H\frac{\rho(d)}d+r_d,
\qquad
|r_d|\leq\rho(d).
\end{equation}

\subsection{Application of the fundamental lemma}

The constants in the fundamental lemma must be uniform in $m$ and $a$.
Define the nonnegative weights
  \[
  a_n=\#\{t\in T_{a,m}:G_{a,m}(t)=n\},
  \qquad X_{a,m}=\sum_n a_n=H.
  \]
Thus, repeated values of $G_{a,m}$ are counted with their multiplicities.
For $d$ supported on $\mathcal P_m$, equation \eqref{eq:localcount} says
\[
\sum_{\substack{n\geq1\\d\mid n}}a_n
=H\frac{\rho(d)}d+r_d.
\]
This is precisely Axiom~1 of \cite[pp.~185--186]{Koukoulopoulos}, with local
density $\nu(d)=\rho(d)$ and remainder $r_d$.  We also have
$0\leq\rho(\ell)<\ell$ for every $\ell\in\mathcal P_m$.

Next we need the dimension condition.  Mertens's estimate and
\eqref{eq:chi} give, uniformly for $w\geq2$,
  \begin{equation}\label{eq:weightedrho}
  \sum_{\substack{\ell\leq w\\ \ell\nmid10}}
  \frac{\rho(\ell)\log\ell}{\ell}
  =\sum_{\ell\leq w}\frac{\log\ell}{\ell}
   +\sum_{\ell\leq w}\chi_5(\ell)\frac{\log\ell}{\ell}
   +O(1)
  =\log w+O(1).
  \end{equation}
Thus, every finite truncation of the base prime set $\mathcal P_0$ satisfies
Axiom~$2'$ of \cite[p.~188]{Koukoulopoulos} with $\kappa=1$, uniformly in
the truncation point.  The additional local condition in that axiom holds,
for example, with $k=2$ and $\varepsilon=1/2$, since
$\rho(\ell)\in\{0,2\}$ and every prime in $\mathcal P_0$ for which
$\rho(\ell)=2$ is at least $11$.  The cited source therefore gives Axiom~2
for every finite truncation of $\mathcal P_0$, with the same absolute
constant $C$.

The deletion of primes dividing $m$ requires a uniformity check.  We do not
claim that Axiom~$2'$ remains unchanged after deleting them.  If
$3/2\leq y_1\leq y_2$, then
\[
\prod_{\substack{\ell\in\mathcal P_m\cap(y_1,y_2]}}
\left(1-\frac{\rho(\ell)}\ell\right)^{-1}
\leq
\prod_{\substack{\ell\in\mathcal P_0\cap(y_1,y_2]}}
\left(1-\frac{\rho(\ell)}\ell\right)^{-1}.
\]
After both sides are restricted to any common finite truncation, the left
side is the product that occurs in Axiom~2 for the truncated set
$\mathcal P_m$.  Hence, every such truncation inherits Axiom~2 from the base
truncation with the same absolute constants.  This proves the required
uniformity in $m$ and $a$.

Set $z=H^{1/20}$ and
  \[
  \mathcal Q_{m,z}=\{\ell\in\mathcal P_m:\ell\leq z\},
  \qquad
  P_m(z)=\prod_{\ell\in\mathcal Q_{m,z}}\ell.
  \]
For sufficiently large $L$, we have $z<L$.  Therefore, every prime value
$G_{a,m}(t)>L$ is counted by the sifted set
$\{t\in T_{a,m}:(G_{a,m}(t),P_m(z))=1\}$.  Let $S_{a,m}$ denote the number
of $t\in T_{a,m}$ for which $G_{a,m}(t)$ is a prime exceeding $L$.  If
$\mathcal Q_{m,z}$ is empty, the estimate below is immediate.  Otherwise,
let $y$ be its largest prime.  We apply
\cite[Theorem 18.11(a), p.~190]{Koukoulopoulos} to the finite prime set
$\mathcal Q_{m,z}$ with $u=10$.  In the notation of that theorem, its main
term is
\[
\left(1+O_{\kappa,C}(u^{-u/2})\right)H V_m(z),
\]
and its remainder is bounded by
\[
O\left(\sum_{\substack{d\leq y^u\\d\mid P_m(z)}}|r_d|\right).
\]
Since $y\leq z$, we have $y^u\leq z^{10}=H^{1/2}$.  Taking an upper bound
and enlarging the remainder sum, we obtain
\begin{equation}\label{eq:Sam}
S_{a,m}
\ll
  HV_m(z)
  +\sum_{\substack{d\leq H^{1/2}\\ d\mid P_m(z)}}|r_d|
  \leq HV_m(z)+\sum_{d\leq H^{1/2}}\rho(d),
\end{equation}
where
\[
V_m(z)=
\prod_{\substack{\ell\leq z\\ \ell\nmid10m}}
\left(1-\frac{\rho(\ell)}\ell\right).
\]
The remainder is easy to control because $\rho(d)\leq\tau(d)$.  Hence,
\[
\sum_{d\leq H^{1/2}}\rho(d)
\leq\sum_{ab\leq H^{1/2}}1
\ll H^{1/2}\log H.
\]
It remains to estimate the sieve product.  First put
  \[
  V_0(z)=\prod_{\substack{\ell\leq z\\ \ell\nmid10}}
  \left(1-\frac{\rho(\ell)}\ell\right).
  \]
By \eqref{eq:split} and the convergence of $\sum_p p^{-2}$, we have
  \[
  \log V_0(z)
  =\sum_{\substack{\ell\leq z\\ \chi_5(\ell)=1}}
  \log\left(1-\frac2\ell\right)
  =-2\sum_{\substack{\ell\leq z\\ \chi_5(\ell)=1}}\frac1\ell+O(1)
  =-\log\log z+O(1).
  \]
The primes dividing $m$ were removed from this product.  To measure the cost
of removing them, define
\[
h(m)=
\prod_{\substack{p\mid m,\ p\neq5\\ \chi_5(p)=1}}
\left(1-\frac2p\right)^{-1}.
\]
Then, omitting from $V_0(z)$ the primes that divide $m$ gives
  \[
  V_m(z)
  =V_0(z)
  \prod_{\substack{\ell\mid m,\ \ell\leq z\\
                    \ell\neq5,\ \chi_5(\ell)=1}}
  \left(1-\frac2\ell\right)^{-1}
  \ll\frac{h(m)}{\log z}
  \ll\frac{h(m)}{\log L}
  \]
uniformly for $m\leq L^\alpha$.  The last comparison follows from
$H\geq L^{1/2}-1$, which gives $\log z\asymp\log L$.  Substituting this into
\eqref{eq:Sam}, we obtain
\begin{equation}\label{eq:Sam2}
S_{a,m}
\ll
\frac{Hh(m)}{\log L}+H^{1/2}\log L.
\end{equation}

\subsection{Summation over complementary factors}

If $i$ is counted by $E_\alpha(L)$,
then for at least one representation $f(i)=mq$, the residue class
$a\equiv i\pmod m$ is one of the $R(m)$ roots and $G_{a,m}(t)=q$ for the
corresponding $t$.  We sum over all such $m$ and all such roots.  This is a
union bound, so an index with more than one representation may be counted
more than once.  Since $m\leq L^\alpha\leq L^{1/2}$, we also have
\[
H\leq\frac Lm+1\leq\frac{2L}{m}.
\]
Put $Y:=L^\alpha$.  Equation \eqref{eq:Sam2} now gives
\begin{equation}\label{eq:Esum}
E_\alpha(L)
\ll
\frac{L}{\log L}
\sum_{m\leq Y}\frac{R(m)h(m)}m
+L^{1/2}\log L
\sum_{m\leq Y}\frac{R(m)}{\sqrt m}.
\end{equation}
The two sums on the right are estimated separately.  For the first, define
$g(m)=R(m)h(m)$.  Equation \eqref{eq:R} and the definition of $h$ show that
$g$ is nonnegative and multiplicative, and that it vanishes outside the
admissible support.  If $p\neq5$ and $\chi_5(p)=1$, then
  \[
  1+\sum_{e\geq1}\frac{g(p^e)}{p^e}
  =1+\frac{2p}{(p-1)(p-2)}
  =1+\frac2p+O(p^{-2}).
  \]
The prime $5$ contributes the bounded local factor $1+1/5$, while inert
primes and $2$ contribute $1$.  Therefore, positivity, the Euler product,
and \eqref{eq:split} give, for $Y\geq1$,
  \begin{align*}
  \sum_{m\leq Y}\frac{R(m)h(m)}m
  &\leq
  \prod_{p\leq Y}\left(1+\sum_{e\geq1}\frac{g(p^e)}{p^e}\right)\\
  &\ll
  \exp\left(
  2\sum_{\substack{p\leq Y\\ \chi_5(p)=1}}\frac1p+O(1)
  \right)
  \ll\log(2Y).
  \end{align*}
For the second sum, we use $R(m)\leq\tau(m)$.  Thus,
\begin{align*}
\sum_{m\leq Y}\frac{R(m)}{\sqrt m}
&\leq
\sum_{ab\leq Y}\frac1{\sqrt{ab}}\\
&\ll
\sum_{a\leq Y}\frac1{\sqrt a}\sqrt{\frac Ya}
\ll\sqrt Y\log(2Y).
\end{align*}
Finally,
  \[
  \log(2Y)\leq\alpha\log L+O(1),
  \]
uniformly for $0<\alpha\leq1/2$.  Substitution of these bounds into
\eqref{eq:Esum}.  This gives
  \[
  E_\alpha(L)
  \ll \alpha L+\frac{L}{\log L}
  +L^{(1+\alpha)/2}(\log L)^2,
  \]
with an absolute implied constant and an $L$-threshold independent of
$\alpha$.  This proves the lemma.
\end{proof}

\section{A low-weight rainbow cover from quadratic congruences}

Lemmas~\ref{lem:largeprime} and \ref{lem:smallcofactor} provide the required
arithmetic input.  The following proposition extracts a family of disjoint
pairs with distinct associated primes and covers the remaining indices
individually.

\begin{proposition}\label{prop:quadratic-cover}
Let $C_0$ be as in Lemma~\ref{lem:smallcofactor}, and set
\[
\alpha_0=\frac1{1000C_0},
\qquad
\delta_0=\frac{\alpha_0}{200},
\qquad
\eta_0=\frac{\delta_0}{2}.
\]
For every sufficiently large $L$, there exist a set
$I\subseteq\{2,\ldots,L-2\}$, distinct primes $q_i$ for $i\in I$, and
distinct primes $r_k$, indexed by
\[
U=\{2,\ldots,L\}\setminus\bigcup_{i\in I}\{i,i+2\},
\]
Put
\[
M=\prod_{i\in I}q_i\prod_{k\in U}r_k.
\]
These objects satisfy the following conditions:
\begin{enumerate}
\item The pairs $\{i,i+2\}$, $i\in I$, are pairwise disjoint, and
$|I|\geq L/100$.
\item For every $i\in I$,
\[
L<q_i<2L^{2-\alpha_0},
\qquad
q_i\mid f(i).
\]
\item For every $k\in U$,
\[
L<r_k\leq20L\log L.
\]
\item The primes $q_i$, $i\in I$, and $r_k$, $k\in U$, are pairwise
distinct.
\item The modulus $M$ satisfies
\[
\log M\leq(1-\delta_0)L\log L,
\qquad
M\leq ((L+1)!)^{1-\eta_0}.
\]
\end{enumerate}
Moreover, every prime divisor of $M$ is less than $L!$.
\end{proposition}

\begin{proof}
The choice of $\alpha_0$ in the proposition and
Lemma~\ref{lem:smallcofactor} give
\[
C_0\alpha_0L=\frac L{1000}.
\]
Since $\alpha_0<1$ is fixed, we also have
\[
C_0\left(\frac L{\log L}
+L^{(1+\alpha_0)/2}(\log L)^2\right)=o(L).
\]
Consequently,
\[
E_{\alpha_0}(L)\leq\frac L{500}
\]
for sufficiently large $L$.  On the other hand,
Lemma~\ref{lem:largeprime} supplies at least $L/5$ indices $i\leq L$ for
which
\[
q_i:=P^+(f(i))>\frac18i\log i.
\]
We discard the indices $i<L/10$, the two endpoints $L-1,L$, and those
counted by $E_{\alpha_0}(L)$.  The number remaining is at least
\[
\frac L5-\frac L{10}-\frac L{500}-2
=\frac{49L}{500}-2
\geq\frac{9L}{100}
\]
for sufficiently large $L$.  For every remaining index,
\[
q_i>\frac18 i\log i
\geq \frac{L}{80}\log\left(\frac{L}{10}\right)>L
\]
once $L$ is sufficiently large.  Write $f(i)=m_iq_i$.  Then
$m_i>L^{\alpha_0}$, since otherwise this index would be counted by
$E_{\alpha_0}(L)$.  Since $f(i)<2L^2$ for $i\leq L$, it follows that
\[
q_i<2L^{2-\alpha_0}.
\]
The identity \eqref{eq:keyidentity} now gives
\[
i!\equiv(i+2)!\pmod{q_i}.
\]

For each remaining index $i$, let
\[
F_i=(q_i,i!,B_i),
\qquad B_i=\{i,i+2\}.
\]
The identity above shows that $B_i\subseteq E_{q_i,i!}$.  We apply
Proposition~\ref{prop:conflict-selection}.  Two designated pairs
$\{i,i+2\}$ and $\{j,j+2\}$ intersect only when $j=i$ or $j=i\pm2$.
Thus, a fixed vertex of the auxiliary conflict graph has at most two
neighbors arising from an overlap of designated pairs.

A fixed prime $q>L$ occurs for at most two starting indices, because all
starts are less than $q$ and $f$ has at most two roots modulo $q$.  Hence a
fixed vertex has at most one additional neighbor of the same prime color.
The conflict graph therefore has maximum degree at most $3$.
Proposition~\ref{prop:conflict-selection}, with unit weights, gives at least
\[
\frac14\cdot\frac{9L}{100}=\frac{9L}{400}
\]
vertices whose associated edges have distinct colors and whose designated
pairs are pairwise disjoint.  For
sufficiently large $L$, we may therefore choose a subfamily indexed by a set
$I$ with
\[
|I|=\left\lceil\frac L{100}\right\rceil.
\]

The set of remaining indices has cardinality $|U|=L-1-2|I|$.  The prime
number theorem gives
\[
\pi(20L\log L)-\pi(L)>2L
\]
for sufficiently large $L$.  After excluding the $|I|$ primes already used
for pairs, more than $2L-|I|>|U|$ primes remain.  Therefore, we may assign a
distinct prime $r_k\in(L,20L\log L]$ to every $k\in U$.

It remains to estimate the product of the selected primes.  Put $r=|I|$ and
$s=|U|=L-1-2r$.  Then
\begin{align*}
\log M
&\leq
r\bigl((2-\alpha_0)\log L+\log2\bigr)
+s\bigl(\log L+\log\log L+\log20\bigr)\\
&=(L-1-\alpha_0r)\log L+s\log\log L+O(L)\\
&\leq
\left(L-1-\frac{\alpha_0L}{100}\right)\log L
+L\log\log L+O(L).
\end{align*}
For sufficiently large $L$, the lower order terms satisfy
\[
L\log\log L+O(L)
\leq\frac{\alpha_0}{200}L\log L,
\]
and therefore
\[
\log M
\leq
\left(1-\frac{\alpha_0}{200}\right)L\log L.
\]
This is the first estimate in the proposition because
$\delta_0=\alpha_0/200$.
Since
\[
\log((L+1)!)=L\log L-L+O(\log L),
\]
we have $\log((L+1)!)\geq L\log L-2L$ for sufficiently large $L$.  Since
$\eta_0=\delta_0/2$, we obtain
\begin{align*}
(1-\eta_0)\log((L+1)!)-\log M
&\geq
(1-\eta_0)(L\log L-2L)-(1-\delta_0)L\log L\\
&=(\delta_0-\eta_0)L\log L-2(1-\eta_0)L>0
\end{align*}
for sufficiently large $L$.  This proves
$M\leq((L+1)!)^{1-\eta_0}$.  The displayed ranges also show that every
selected prime is less than $L!$ once $L$ is sufficiently large.  Finally,
we take the maximum of the finitely many absolute thresholds imposed above.
This completes the construction.
\end{proof}

\begin{proof}[Proof of Theorem~\ref{thm:covercost}]
Choose the sets and primes from Proposition~\ref{prop:quadratic-cover}.  For
each $i\in I$, the identity \eqref{eq:keyidentity} and the divisibility
$q_i\mid f(i)$ give
\[
i!\equiv(i+2)!\pmod{q_i}.
\]
Hence the edge $e_{q_i,i!}$ covers both $i$ and $i+2$.  For each $k\in U$,
the edge $e_{r_k,k!}$ covers $k$.  All the primes are distinct, so these
edges form a rainbow edge cover $\mathcal C_L$.  Proposition
\ref{prop:quadratic-cover} gives
\[
\kappa(L)\leq w(\mathcal C_L)
=\log M\leq(1-\delta_0)L\log L.
\]
Taking $\delta=\delta_0$ proves the theorem.
\end{proof}

\section{From rainbow covers to rough integers}

The following transference theorem applies to any rainbow edge cover with a
power-saving modulus.

\begin{theorem}[Transference theorem]\label{thm:transference}
Fix $0<\eta<1$.  For every sufficiently large $L$, depending only on
$\eta$, let $\mathcal C$ be a rainbow edge cover of $\mathcal H_L$ and put
\[
X=(L+1)!,
\qquad
M=M(\mathcal C).
\]
If $M\leq X^{1-\eta}$, then at least
\[
\frac{X^\eta}{100\log L}
\]
integers $n\in(2L!,X]$ satisfy $P^-(n)>L$ and have $n-k!$ composite for
every $1\leq k\leq L$.
\end{theorem}

\begin{proof}
Write
\[
\mathcal C=\{e_{q,a_q}:q\in Q\}.
\]
The primes in $Q$ are distinct, so the Chinese remainder theorem gives a
unique residue class $A\pmod M$ satisfying
\begin{equation}\label{eq:A}
A\equiv a_q\pmod q\qquad(q\in Q).
\end{equation}
Every $a_q$ is nonzero, and hence $(A,M)=1$.

Let
\[
\mathcal T=\{t\in\mathbb Z:2L!<A+tM\leq X\},
\qquad
H=|\mathcal T|.
\]
The set $\mathcal T$ is an interval of $H$ consecutive integers.  Its length
gives
\[
H\geq\frac{X-2L!}{M}-1
=\frac{L-1}{L+1}\frac XM-1.
\]
For $L\geq3$, the right side is at least $X/(2M)-1$.  Also,
$M\leq X^{1-\eta}$ gives $X/M\geq X^\eta\to\infty$.  Therefore, for all
sufficiently large $L$,
\begin{equation}\label{eq:Hlower}
H\geq\frac{X}{3M}\geq\frac{X^\eta}{3}.
\end{equation}

Put
\[
W=\prod_{p\leq L}p.
\]
Every prime factor of $M$ exceeds $L$, so $(M,W)=1$.  Writing
$\vartheta(y)=\sum_{p\leq y}\log p$, the prime number theorem gives
\[
\log W=\vartheta(L)=L+o(L).
\]
Stirling's formula and \eqref{eq:Hlower} give
\begin{align*}
\log\frac{H}{2W}
&\geq \eta\log X-\log6-\vartheta(L)\\
&=\eta L\log L-(1+\eta+o(1))L-\log6
\longrightarrow+\infty.
\end{align*}
Thus $H\geq2W$ for sufficiently large $L$.

Multiplication by $M$ is invertible modulo $W$.  Hence the map
\[
t\longmapsto A+tM\pmod W
\]
is a bijection on $\mathbb Z/W\mathbb Z$.  Every block of $W$ consecutive
values of $t$ contains exactly $\varphi(W)$ values for which
$(A+tM,W)=1$.  Since $H\geq2W$, at least
\[
\left\lfloor\frac HW\right\rfloor\varphi(W)
\geq\frac H2\frac{\varphi(W)}W
\]
values $t\in\mathcal T$ have this property.  Mertens's product theorem gives
\[
\frac{\varphi(W)}W
=\prod_{p\leq L}\left(1-\frac1p\right)
\geq\frac1{4\log L}
\]
for sufficiently large $L$.  It follows that the number of integers
$n=A+tM$ in $(2L!,X]$ with $(n,W)=1$ is at least
\[
\frac{H}{8\log L}
\geq\frac{X}{24M\log L}
\geq\frac{X^\eta}{24\log L}
\geq\frac{X^\eta}{100\log L}.
\]
Every such $n$ is $L$-rough.

It remains to check the factorial differences.  Let $2\leq k\leq L$.
Since $\mathcal C$ is a rainbow edge cover, there is a $q\in Q$ such that
$k\in E_{q,a_q}$.  Equations \eqref{eq:A} and the definition of the fiber
give
\[
q\mid n-k!.
\]
This divisor is proper.  Indeed, $q<L!$ by Definition~\ref{def:fibres}, and
\[
n-k!>2L!-k!\geq L!.
\]
Thus $1<q<n-k!$, and $n-k!$ is composite.

Finally, $(n,W)=1$ implies that $n$ is odd.  Since $n>2L!\geq4$, the number
$n-1$ is an even integer greater than $2$.  Hence $n-1$ is composite, which
settles the case $k=1$.
\end{proof}

\begin{proof}[Proof of Theorem~\ref{thm:main}]
Let $\delta\in(0,1)$ be the constant in Theorem~\ref{thm:covercost}, and
put $\eta=\delta/2$.  By Propositions~\ref{prop:baseline} and
\ref{prop:crt-equivalence}, for all sufficiently large $L$ there is a
minimum-weight rainbow edge cover $\mathcal C_L$ with
\[
\log M(\mathcal C_L)=\kappa(L)
\leq(1-\delta)L\log L.
\]
Also, Stirling's formula gives
\[
\log X\geq L\log L-2L
\]
for sufficiently large $L$.  Therefore,
\begin{align*}
(1-\eta)\log X-\log M(\mathcal C_L)
&\geq(\delta-\eta)L\log L-2(1-\eta)L\\
&=\frac{\delta}{2}L\log L-2(1-\eta)L>0
\end{align*}
once $L$ is large enough.  Thus
\[
M(\mathcal C_L)\leq X^{1-\eta}.
\]
Theorem~\ref{thm:transference} now gives the result after increasing the
absolute threshold $L_1$ if necessary.
\end{proof}

\section{Concluding remarks}

The hypergraph $\mathcal H_L$ is not tied to gap two.  By
Proposition~\ref{prop:criterion}, a prime $q>L$ dividing
\[
D_{i,j}=\prod_{r=i+1}^{j}r-1
\]
places $i$ and $j$ in a common edge.  For example, the gap-three identity
\[
\frac{(x+1)!}{(x-2)!}=x^3-x
\]
gives a common edge containing $x-2$ and $x+1$ whenever
$q\mid x^3-x-1$.  The corresponding auxiliary conflict graph has bounded
degree when $L<q\leq L^{3/2}$, so the independent-set lemma selects a fixed
proportion of any available family with distinct colors and disjoint
designated pairs.  The unresolved arithmetic question is whether such
low-weight prime divisors occur for a positive proportion of the indices.

\section*{Acknowledgments}

This work was carried out as part of the 2026 Research in Industrial Projects
for Students program at the Institute for Pure \& Applied Mathematics,
UCLA.  The author thanks IPAM for providing a collaborative research
environment and acknowledges OpenAI, the industrial sponsor of the project,
for its support and for providing access to the computational tools used
during the project.

\section*{Declaration of generative AI and AI-assisted technologies in the
manuscript preparation process}

During the research and preparation of this work, the author used OpenAI's
ChatGPT and Codex for exploratory work, computation, literature searches,
algebraic checks, proof organization, and editorial revision.  The author
reviewed and verified the resulting material, edited the manuscript as
needed, and takes full responsibility for the content of the article.
\begin{thebibliography}{99}

\bibitem{BanksLucaShparlinskiStichtenoth}
W.~D. Banks, F.~Luca, I.~E. Shparlinski, and H.~Stichtenoth,
\emph{On the value set of $n!$ modulo a prime},
Turkish J. Math. \textbf{29} (2005), no.~2, 169--174.

\bibitem{Bloom}
T.~F. Bloom,
\emph{Erd\H{o}s Problem 1059},
\href{https://www.erdosproblems.com/1059}
{erdosproblems.com/1059}, accessed August 9, 2026.

\bibitem{GaraevLucaShparlinski}
M.~Z. Garaev, F.~Luca, and I.~E. Shparlinski,
\emph{Exponential sums and congruences with factorials},
J. Reine Angew. Math. \textbf{584} (2005), 29--44.

\bibitem{Guy}
R.~K. Guy,
\emph{Unsolved Problems in Number Theory},
3rd ed., Problem Books in Mathematics, Springer, New York, 2004,
Problem A2, pp.~10--12.

\bibitem{Hu}
X.~Hu,
\emph{Factorial residues modulo a prime: beyond the square-root bound},
\href{https://arxiv.org/abs/2608.01781}{arXiv:2608.01781} (2026).

\bibitem{KlurmanMunsch}
O.~Klurman and M.~Munsch,
\emph{Distribution of factorials modulo $p$},
J. Th\'eor. Nombres Bordeaux \textbf{29} (2017), no.~1, 169--177.

\bibitem{Koukoulopoulos}
D.~Koukoulopoulos,
\emph{The Distribution of Prime Numbers},
Graduate Studies in Mathematics, vol.~203,
American Mathematical Society, Providence, RI, 2019.

\end{thebibliography}

\end{document}
